\documentclass{article}
\usepackage{amsmath, amssymb, amsthm}
\usepackage{hyperref}
\usepackage{geometry}
\geometry{margin=1in}

\title{Erd\H{o}s Problem \#740}
\date{Last updated: 2025-08-31}
\author{Source: erdosproblems.com}

\begin{document}
\maketitle

\noindent\textbf{Status:} OPEN \quad \textbf{No prize}

\noindent\textbf{Tags:} graph theory, chromatic number

\medskip
\noindent\textbf{Problem Statement:}

Let $\mathfrak{m}$ be an infinite cardinal and $G$ be a graph with chromatic number $\mathfrak{m}$. Let $r\geq 1$. Must $G$ contain a subgraph of chromatic number $\mathfrak{m}$ which does not contain any odd cycle of length $\leq r$?

\bigskip
\noindent\textbf{Remarks:}

A question of Erd\H{o}s and Hajnal. R\"{o}dl proved this is true if $\mathfrak{m}=\aleph_0$ and $r=3$ (see [108] for the finitary version).

More generally, Erd\H{o}s and Hajnal asked must there exist (for every cardinal $\mathfrak{m}$ and integer $r$) some $f_r(\mathfrak{m})$ such that every graph with chromatic number $\geq f_r(\mathfrak{m})$ contains a subgraph with chromatic number $\mathfrak{m}$ with no odd cycle of length $\leq r$?

Erd\H{o}s \cite{Er95d} claimed that even the $r=3$ case of this is open: must every graph with sufficiently large chromatic number contain a triangle free graph with chromatic number $\mathfrak{m}$?

In \cite{Er81} Erd\H{o}s also asks the same question but with girth (i.e. the subgraph does not contain any cycle at all of length $\leq C$).

\bigskip
\noindent\textbf{References:}

[Er81] Erd\H{o}s, P., On the combinatorial problems which I would most like to see solved. Combinatorica (1981), 25-42.
 
 [Er95d] Erd\H{o}s, Paul, On some problems in combinatorial set theory. Publ. Inst. Math. (Beograd) (N.S.) (1995), 61-65.

\bigskip
\noindent\small{Source: \url{https://www.erdosproblems.com/740}}
\end{document}
